\section{Background and Related Work}
\label{sec:related}
\label{sec:background}

\paragraph{Sharded BFT and quorum intersection.}
We assume the now-standard committee-based design~\cite{luu2016elastico,kokoris2018omniledger,zamani2018rapidchain}
in which each shard runs a quorum-based BFT protocol~\cite{castro1999pbft,yin2019hotstuff,buchman2018tendermint}.
The safety of a single shard rests on quorum intersection (Lemma~\ref{lem:intersection}). Cross-shard
atomicity is typically achieved with a two-phase-commit-style protocol coordinated either by the
client or by the shards~\cite{kokoris2018omniledger}; both variants are known to be fragile, as we
discuss below.

\paragraph{Single-shard compromise and adaptive corruption.}
That a committee exceeding its Byzantine threshold compromises all transactions touching it is
identified as an outstanding problem in the canonical systematization~\cite{bano2019sok}. A long
line of sharding work treats the Byzantine fraction as \emph{exogenous} and asks how to make a
captured committee less likely---large committees and unbiasable assignment~\cite{luu2016elastico},
periodic reconfiguration and bounded-adversary epoch models~\cite{zamani2018rapidchain}---or how to
detect capture after the fact. We instead ask what it \emph{costs} to produce the captured committee
once membership is revealed, and what the captor can then extract before accountability resolves.

\paragraph{Adaptive corruption and corrupted-shard tolerance.}
Closest in setting are adaptive-corruption analyses of sharding. \cite{avarikioti2019divide}
formalizes when a sharded ledger stays consistent and scalable against epoch-adaptive versus fully
adaptive adversaries (requiring succinct state-update proofs at epoch boundaries), and
\cite{rana2022free2shard} defends a globally sub-threshold but locally concentrated adaptive adversary
by letting honest nodes re-allocate across shards; \cite{david2022gearbox} minimizes committee size by
separating safety from liveness, which only \emph{amplifies} the leverage of our attack, since smaller
committees are cheaper to bribe. Most directly, the concurrent Camael design~\cite{liu2025camael}
formalizes exactly our structural premise---a single shard exceeding its local BFT threshold (up to
$2/3$ corrupt) while global stake stays below $1/3$---and builds a \emph{defense} that detects
concealed safety violations via snapshots, convicts offenders, and reconfigures the shard. These works
treat the captured committee as a \emph{corruption}-budget event and target consistency, liveness, or
remediation; we treat it as an \emph{economic procurement} event and ask what the captor extracts
before accountability resolves. Tellingly, Camael's snapshot-and-conviction pipeline is itself an
instance of the accountability latency $\Tacc$ our race turns on.

\paragraph{Bribery and crypto-economic security.}
The economic study of buying consensus predates trustless contracts: \cite{bonneau2016rent} showed a
\emph{rented} majority has no long-term stake to protect and is cheap to mobilize, and the
$P{+}\epsilon$ argument~\cite{buterin2015pepsilon} showed a credible conditional payment can make the
adversary's preferred action dominant at near-zero realized cost. Bribing consensus participants via
smart contracts originates with~\cite{mccorry2018smart};
\cite{winzer2019temporary} studies temporary censorship by rational miners (with bribes that are
\emph{not} trustless, as the bribee must trust the briber). \cite{tran2023crosschain} proposes
cross-chain bribery contracts that bribe rational miners in a target network at low cost.
\cite{judmayer2021sok,judmayer2021paytowin} systematize these \emph{algorithmic
incentive-manipulation} attacks and demonstrate cheap, crowdfundable, cross-chain bribery funding a
double-spend on a target chain, and out-of-band collusion contracts~\cite{zhang2023outofband} make
defection the strict equilibrium for a sub-threshold rational coalition, even sketching evasion of
slashing by censoring the proof; that proposer--user side contracts are a fundamental incentive limit
is established by the impossibility of side-contract-proof transaction-fee
mechanisms~\cite{roughgarden2021tfm,chung2023foundations}.
\cite{karakostas2024bribing} gives a game-theoretic treatment of bribery aimed at PoS \emph{safety},
distinguishing \emph{guided} bribes (paid for prescribed behavior) from \emph{effective} bribes
(paid conditional on attack success), and shows that counterincentives---specifically slashing and
dilution---characterize when the attack becomes unprofitable. Most recently,
\cite{sookitoth2025bribers} implements and evaluates three trustless bribery mechanisms on a deployed
PoS system---buying votes to fork, inducing validator exits, and auctioning randomness-beacon
(RANDAO) manipulation---and explicitly debunks the belief that crypto-economic security equals the
total staked value. Modern variants coerce attesters through credible
commitments~\cite{sarenche2024commitment} or buy block delay for MEV without triggering slashing at
all~\cite{yang2024bride}, while \cite{newman2023decentralization} shows the equilibrium per-validator
bribe can fall well below naive expected-loss bounds via low pivotality. Related work studies smart
collusion and whistleblowing as a counterforce to
such contracts~\cite{kelkar2025omerta}, the private liquefaction of staked assets to enable
vote-buying~\cite{austgen2025liquefaction}, and grinding/manipulation of blockchain randomness
beacons~\cite{gazi2025grinding}. These works establish that buying equivocation is feasible,
trustless, and far cheaper than the total-stake figure suggests. They do \emph{not} model the
small-committee quorum structure, the cross-shard realization of the resulting violation, or the way
accountability latency \emph{prices} the bribe.

\paragraph{Economic limits, slashable safety, and accountability latency.}
A separate line bounds attacks economically rather than by a Byzantine fraction. Budish's
flow-versus-stock argument~\cite{budish2018economic} and its permissionless-consensus
generalization~\cite{budish2024economic} establish that an attack is deterred only when its cost
exceeds its profit and---decisively for us---that under partial synchrony above the $1/3$ threshold
targeted punishment \emph{cannot} be enforced before the violator liquidates, an intrinsic
latency-bounded limit on accountability. STAKESURE makes the cost-of-corruption versus
profit-from-corruption gate explicit and caps attacker profit by gating off-chain and bridge value
release on a reversion window~\cite{deb2024stakesure}; CoBRA bounds finalizable value below the
slashable stake so a rational validator's gain cannot exceed its bond~\cite{avarikioti2025cobra}; and
Bitcoin-enhanced PoS makes the unbonding delay the operative security variable, racing slashing
against withdrawal and gating safety on external finality~\cite{tas2023bitcoin}. These results rest on
an accountability substrate we also assume: identifiability of culprits after a
fork~\cite{civit2021polygraph}, per-protocol forensic support fixing how many of the necessary
equivocators are provably catchable~\cite{sheng2021forensics}, the accountable-finalized prefix
underlying finality-gating~\cite{neu2022accountability}, and the detection-and-recovery machinery
whose latency we abstract as $\Tacc$~\cite{gong2025recover}. All study a monolithic, single-chain, or
restaking~\cite{durvasula2025restaking} setting in which the adversary owns stake or value is
throttled, and most assume slashing is effectively immediate. We instead study a globally
\emph{sub-threshold} adversary who, after committee assignment, \emph{buys} the $F{+}1$ equivocators
one $N{=}3F{+}1$ committee makes necessary, prices that bribe by the probability $q$ that the proof is
enforced before stake becomes unslashable, and realizes value across the shard/bridge boundary in
$\Tsettle$ before $\Tacc$; our finality-gating immunity (Proposition~\ref{prop:immunity}) specializes
the accountable-finalized-prefix idea to cross-shard acceptance.

\paragraph{Bribing against an accountability window.}
Closest to our race in spirit is the economic-censorship game of~\cite{berger2025censorship}, in which
an attacker bribes block proposers to censor a defender's challenge transactions and \emph{exhaust the
challenge-period time budget}, with the attacker/defender budget ratio deciding who wins; the
analogous failure where winning a fraud proof is itself unprofitable is studied
by~\cite{lee2025hollow}. This is the same move our $\Tsettle<\Tacc$ gate captures---paying to outlast
accountability---but their adversary censors a single defender's transaction in an
L1$\rightarrow$L2 challenge window, whereas ours buys the quorum-intersection equivocators of an
$N{=}3F{+}1$ committee and races \emph{irreversible cross-shard/bridge} settlement against
detection-and-slashing rather than depleting a fixed challenge clock. The dual defensive
primitive---a remote unbonding protocol that guarantees slashing \emph{before} stake is
withdrawn---is exactly the ``slash before unslashable'' condition our $q$
measures~\cite{dong2024remotestaking}, as is the ``correct until withdrawal'' collateral with
committee resets of~\cite{baron2025aegis}.

\paragraph{Cross-shard, cross-chain, and settlement-before-finality value extraction.}
\cite{sonnino2020replay} presents replay attacks on cross-shard consensus that enable double-spends
and even minting \emph{without subverting any node} and even when Byzantine safety holds---a
different vector from ours (protocol-level message replay, not committee corruption), defended by
their Byzcuit protocol, but clear evidence that cross-shard settlement is a realizable value sink. In
the cross-chain setting, extracting value on a destination and then reverting the source---the
canonical bridge double-mint and reorg-double-spend pattern~\cite{zamyatin2021sok,lee2022sokbridge}---is
precisely the timing move our $\Tsettle<\Tacc$ gate formalizes, but there the reversal is an
opportunistic reorg against a bridge that acts before finality~\cite{schwarzschilling2022three}, not a
reversal manufactured \emph{inside} the protocol by a bribed committee. Miner/maximal-extractable-value
research~\cite{daian2020flashboys} further shows that consensus-level manipulation has directly
monetizable value, and empirical and cross-domain studies quantify the magnitude of such extractable
value~\cite{qin2022quantifying,mcmenamin2023sok}; together these motivate treating $V$ as bounded by
reachable liquidity rather than local balances.

\paragraph{Accountable BFT, slashing, and reconfiguration.}
Accountable consensus turns equivocation into a transferable fraud proof and slashes the offending
stake~\cite{buterin2017casper,buchman2018tendermint}; this is the defender we assume, and the reason
the bare composition of ``small committee'' and ``buyable faults'' is not yet an attack. Sharded
systems additionally rotate committees across epochs and hand off shard state and cross-shard
artifacts at the boundary~\cite{zamani2018rapidchain,kokoris2018omniledger}. Prior work studies
reconfiguration as a means of \emph{raising} adversarial cost (a fresh committee each epoch); we
study how the boundary interacts with accountability, and identify exactly when handoff relocates the
race rather than closing it (\S\ref{sec:reconf}).

\paragraph{Positioning.}
Table~\ref{tab:positioning} situates this work against the closest prior art across the ten
capability dimensions that jointly define the attack. The five \emph{load-bearing} axes are
whether the work \emph{prices} the corruption, whether it is \emph{committee-local}, whether it
models \emph{cross-shard settlement}, whether it models the \emph{accountability race} (value
before slashing), and whether it treats \emph{reconfiguration handoff}; we add five finer
capability columns---a \emph{trustless} payment channel, a \emph{safety} (rather than liveness or
censorship) target, a bribe priced by slashing \emph{latency} $q$, an \emph{immunity} condition,
and an \emph{implemented/measured} artifact---so the comparison is granular rather than coarse.
Read by rows, every prior mechanism leaves at least four columns empty; read by the bottom block,
no prior work addresses more than a handful of the dimensions. Only this work spans all of them,
and that intersection---plus the explicit viability and immunity conditions---is the contribution
of this paper.

\begin{table*}[t]
\centering
\caption{Positioning against the closest prior work across the capability dimensions of the attack.
\yes~addresses it, \partialmark~partially/implicitly, \nomark~does not, \namark~not applicable to
that work's setting. Columns: bribe \emph{priced}; \emph{trustless} payment (no briber--bribee
trust); targets \emph{safety} (vs.\ liveness/censorship); \emph{committee-local} sub-threshold
quorum; cross-shard/bridge \emph{value} realization; the \emph{accountability race} (value before
slashing); bribe priced by slashing \emph{latency} $q$; an \emph{immunity}-by-construction
condition; \emph{reconfiguration} handoff; and an \emph{implemented/measured} artifact. This work
is the only row populated across every dimension.}
\label{tab:positioning}
\scriptsize
\renewcommand{\arraystretch}{1.12}
\setlength{\tabcolsep}{4.5pt}
\begin{tabular}{@{}l l *{10}{c}@{}}
\toprule
\textbf{Work} & \textbf{Setting}
 & \rotatebox[origin=l]{90}{Bribe \emph{priced}\,}
 & \rotatebox[origin=l]{90}{Trustless\,}
 & \rotatebox[origin=l]{90}{Targets \emph{safety}\,}
 & \rotatebox[origin=l]{90}{Committee-local\,}
 & \rotatebox[origin=l]{90}{Cross-shard $V$\,}
 & \rotatebox[origin=l]{90}{Accountability race\,}
 & \rotatebox[origin=l]{90}{Latency-priced $q$\,}
 & \rotatebox[origin=l]{90}{Immunity condition\,}
 & \rotatebox[origin=l]{90}{Reconfig.\ handoff\,}
 & \rotatebox[origin=l]{90}{Impl./measured\,} \\
\midrule
\textbf{This work} & sharded
 & \yes & \yes & \yes & \yes & \yes & \yes & \yes & \yes & \yes & \yes \\
\midrule
\multicolumn{12}{@{}l}{\textit{Bribery mechanisms}}\\
Bribing miners~\cite{mccorry2018smart} & monolithic
 & \yes & \yes & \partialmark & \nomark & \nomark & \nomark & \nomark & \nomark & \nomark & \yes \\
Cross-chain bribery~\cite{tran2023crosschain} & cross-chain
 & \yes & \yes & \partialmark & \nomark & \partialmark & \nomark & \nomark & \nomark & \nomark & \partialmark \\
Pay-to-win~\cite{judmayer2021paytowin} & cross-chain
 & \yes & \yes & \partialmark & \nomark & \partialmark & \nomark & \nomark & \nomark & \nomark & \partialmark \\
Out-of-band collusion~\cite{zhang2023outofband} & monolithic
 & \yes & \yes & \nomark & \partialmark & \nomark & \partialmark & \nomark & \nomark & \nomark & \nomark \\
Bribing \& counterincentives~\cite{karakostas2024bribing} & monolithic
 & \yes & \namark & \yes & \nomark & \partialmark & \partialmark & \nomark & \nomark & \nomark & \nomark \\
BriDe Arbitrager~\cite{yang2024bride} & monolithic
 & \yes & \yes & \nomark & \partialmark & \nomark & \nomark & \nomark & \nomark & \nomark & \yes \\
Trustless bribery contracts~\cite{sookitoth2025bribers} & monolithic
 & \yes & \yes & \partialmark & \nomark & \nomark & \nomark & \nomark & \nomark & \nomark & \yes \\
\addlinespace[1pt]
\multicolumn{12}{@{}l}{\textit{Economic limits and slashable safety}}\\
Economic limits~\cite{budish2024economic} & monolithic
 & \namark & \namark & \yes & \nomark & \nomark & \partialmark & \partialmark & \nomark & \nomark & \nomark \\
STAKESURE~\cite{deb2024stakesure} & restaking
 & \partialmark & \namark & \partialmark & \nomark & \partialmark & \partialmark & \nomark & \yes & \nomark & \nomark \\
CoBRA~\cite{avarikioti2025cobra} & monolithic
 & \namark & \namark & \yes & \nomark & \partialmark & \partialmark & \nomark & \yes & \nomark & \partialmark \\
Bitcoin-enhanced PoS~\cite{tas2023bitcoin} & monolithic
 & \namark & \namark & \yes & \nomark & \nomark & \partialmark & \partialmark & \yes & \partialmark & \partialmark \\
Remote staking~\cite{dong2024remotestaking} & restaking
 & \namark & \namark & \yes & \nomark & \partialmark & \yes & \yes & \partialmark & \partialmark & \nomark \\
\addlinespace[1pt]
\multicolumn{12}{@{}l}{\textit{Accountability-window and sharding}}\\
Economic censorship games~\cite{berger2025censorship} & rollup
 & \partialmark & \namark & \nomark & \nomark & \nomark & \yes & \partialmark & \nomark & \nomark & \nomark \\
Corrupted-shard tolerance~\cite{liu2025camael} & sharded
 & \namark & \namark & \yes & \yes & \partialmark & \nomark & \nomark & \nomark & \partialmark & \yes \\
Cross-shard replay~\cite{sonnino2020replay} & sharded
 & \namark & \namark & \yes & \yes & \yes & \nomark & \nomark & \nomark & \nomark & \yes \\
Free2Shard~\cite{rana2022free2shard} & sharded
 & \namark & \namark & \yes & \yes & \nomark & \nomark & \nomark & \nomark & \partialmark & \partialmark \\
\bottomrule
\end{tabular}
\end{table*}
